\capredactadotrue
\section{Análisis Estratégico}\label{sec:analisis-estrategico}

El análisis procede por aproximaciones sucesivas: parte del contexto general en que se lanzará el proyecto (local, económico, político y tecnológico), lo traduce en un escenario-meta y en oportunidades y amenazas sectoriales, y desciende luego al análisis de la competencia y de los factores que determinan el éxito del negocio.

\subsection{Análisis de contexto}

\subsubsection{Escenario local}\label{sec:escenario-local}
El proyecto se lanzará inicialmente en Argentina, comenzando por la ciudad de Rosario, lo que permite aprovechar el conocimiento directo del mercado local y la cercanía con los potenciales usuarios.

En los últimos años se observa un crecimiento sostenido de las comunidades de juegos de mesa y de TTRPG en ciudades como Rosario, Córdoba y Buenos Aires. Asimismo, la expansión de las redes sociales, de Discord y de las plataformas de \textit{streaming} facilitó la difusión del pasatiempo y la incorporación de nuevos jugadores.

El carácter digital del producto permite una rápida escalabilidad hacia otros mercados latinoamericanos e internacionales sin necesidad de infraestructura física significativa.

\subsubsection{Factores económicos}\label{sec:economicos}
El contexto económico presenta tanto desafíos como oportunidades para el proyecto.

Por un lado, la inflación, las fluctuaciones del tipo de cambio y la incertidumbre económica pueden afectar la capacidad de gasto de los consumidores en productos vinculados al entretenimiento y a los pasatiempos. Esto puede reducir la disposición de algunos usuarios a contratar servicios premium o a realizar compras dentro de la plataforma.

Por otro lado, el proyecto ofrece una alternativa digital de bajo costo frente a otros gastos asociados al pasatiempo, como la adquisición de manuales importados, miniaturas, mapas físicos o suscripciones a plataformas internacionales, cuyos precios suelen estar dolarizados. Además, el modelo SaaS permite una estructura de costos flexible y escalable, que reduce la necesidad de inversiones significativas en infraestructura física. Por último, la posibilidad de acceder a mercados internacionales disminuye la dependencia exclusiva de la situación económica local y diversifica las fuentes de ingresos.

\subsubsection{Factores políticos}\label{sec:politicos}
El entorno político e institucional favorece el crecimiento de los servicios digitales mediante iniciativas de transformación digital, la expansión de la conectividad y la adopción de tecnologías basadas en Internet.

La actividad del proyecto se encuentra alcanzada por normativas sobre comercio electrónico, protección de datos personales, propiedad intelectual y contratación de servicios digitales. En particular, resulta relevante garantizar el tratamiento seguro de la información de los usuarios y una adecuada gestión del contenido generado por la comunidad. Además, posibles modificaciones en las regulaciones fiscales sobre servicios digitales, medios de pago electrónicos o comercialización internacional pueden influir en el modelo de negocios y en la rentabilidad futura.

La naturaleza digital de la plataforma permite operar con relativa independencia de las barreras geográficas y regulatorias locales, aunque será necesario monitorear de forma continua los cambios normativos, tanto en Argentina como en los potenciales mercados internacionales.

\subsubsection{Factores tecnológicos}\label{sec:tecnologicos}
Los factores tecnológicos favorecen el desarrollo del proyecto por los siguientes motivos:
\begin{itemize}
  \item amplia disponibilidad de servicios \textit{cloud} escalables;
  \item crecimiento de las arquitecturas SaaS;
  \item adopción masiva de dispositivos móviles;
  \item acceso generalizado a Internet;
  \item madurez de las tecnologías modernas de desarrollo web.
\end{itemize}
Asimismo, la incorporación futura de herramientas de inteligencia artificial podría facilitar la generación de contenido narrativo, los resúmenes automáticos de campañas y la asistencia al DM.

\subsubsection{Escenario-meta}
Se prevé un escenario favorable, caracterizado por:
\begin{itemize}
  \item crecimiento continuo del mercado de TTRPG;
  \item mayor digitalización de las actividades recreativas;
  \item incremento del uso de herramientas colaborativas en línea;
  \item desarrollo de comunidades globales conectadas.
\end{itemize}
En este escenario, la plataforma busca posicionarse como la principal herramienta de gestión de campañas presenciales, alcanzar una masa crítica de usuarios y consolidar un ecosistema de contenido generado por la comunidad.

\subsubsection{Análisis sectorial: oportunidades y amenazas}\label{sec:oportunidades-amenazas}
\paragraph{Oportunidades.}
\begin{itemize}
  \item Crecimiento global del mercado de TTRPG.
  \item Incremento de la cantidad de jugadores jóvenes.
  \item Falta de herramientas especializadas para campañas presenciales.
  \item Expansión de los modelos SaaS.
  \item Desarrollo de comunidades digitales.
  \item Potencial de internacionalización.
\end{itemize}

\paragraph{Amenazas.}
\begin{itemize}
  \item Competidores establecidos, como Roll20, Foundry VTT y D\&D Beyond.
  \item Cambios en las políticas de licencias de propiedad intelectual.
  \item Aparición de nuevos competidores tecnológicos.
  \item Baja disposición inicial a pagar por herramientas digitales dentro del pasatiempo.
  \item Dependencia de plataformas externas para la distribución y la publicidad.
\end{itemize}

\subsection{Análisis de la competencia}

\subsubsection{Principales competidores directos}\label{sec:competidores}
\paragraph{Roll20.}
Es una de las plataformas más reconocidas del ecosistema de juegos de rol de mesa. Su propuesta consiste en permitir partidas completamente en línea, mediante un entorno virtual que replica muchos de los elementos de una mesa física. Integra mapas interactivos, fichas de personaje digitales, sistemas de lanzamiento de dados y herramientas de comunicación, de modo que jugadores ubicados en lugares distintos participan simultáneamente de una misma sesión.

Si bien comparte con el proyecto el objetivo de asistir a jugadores y directores de juego en la gestión de campañas, su enfoque se orienta principalmente a la virtualización completa de la experiencia. La solución propuesta, en cambio, busca complementar las partidas presenciales mediante herramientas organizativas y de acceso a la información, sin reemplazar la dinámica social propia de los juegos de rol de mesa tradicionales.

\paragraph{Foundry Virtual Tabletop.}
Es una plataforma especializada en la creación de entornos virtuales altamente personalizables. Se distingue por su nivel de configuración y extensibilidad: permite incorporar módulos desarrollados por la comunidad para automatizar reglas, gestionar campañas complejas o agregar funcionalidades para distintos sistemas de juego.

Su principal fortaleza es la capacidad de recrear digitalmente una partida mediante mapas dinámicos, efectos visuales y automatización avanzada. Esa misma complejidad puede ser una barrera de entrada para quienes solo necesitan organizar información, planificar sesiones o gestionar campañas presenciales. El proyecto apunta a un segmento diferente, centrado en la simplicidad de uso y en la centralización de la información relevante para la experiencia física de juego.

\paragraph{D\&D Beyond.}
Es la plataforma oficial de contenido digital de \textit{Dungeons \& Dragons} y una de las herramientas más utilizadas por jugadores y \textit{Dungeon Masters} de ese sistema. Su propuesta principal consiste en digitalizar manuales, reglas y hojas de personaje, y permitir el acceso rápido a la información oficial desde cualquier dispositivo.

\paragraph{World Anvil.}
Es una plataforma orientada a la construcción de mundos ficticios, la documentación narrativa y la organización de información compleja. La utilizan tanto directores de juego como escritores y creadores de contenido que desean estructurar historias, personajes, ubicaciones y acontecimientos dentro de universos extensos.

Su fortaleza reside en la profundidad de sus capacidades de documentación y de construcción de mundos, que permiten elaborar verdaderas enciclopedias digitales de campañas y escenarios. Sin embargo, muchas de sus funcionalidades se enfocan en la planificación narrativa y la elaboración de contenido, más que en la gestión operativa de sesiones y personajes durante la partida. El proyecto se diferencia al priorizar herramientas prácticas para el día a día de jugadores y \textit{Dungeon Masters}, que faciliten la consulta y la administración de la información durante el desarrollo de las campañas.

\paragraph{Síntesis.}
En términos generales, las soluciones más consolidadas se concentran en dos enfoques. Por un lado, las plataformas orientadas a la virtualización completa de la experiencia de juego, como Roll20 y Foundry VTT. Por otro, las herramientas enfocadas en la documentación, la gestión de contenido y el acceso a reglas, como D\&D Beyond y World Anvil.

La propuesta desarrollada busca posicionarse en un espacio menos explotado: el apoyo a las campañas presenciales mediante una plataforma híbrida que centralice la información, facilite la coordinación entre los participantes y reduzca las tareas administrativas de gestión de una campaña. Esta diferenciación permite competir no por la sustitución de las partidas físicas, sino por la mejora de la experiencia tradicional de juego, una oportunidad de mercado poco atendida por las principales soluciones existentes.

\subsubsection{Análisis de cadena de valor}\label{sec:cadena-valor}
La cadena de valor de \nombreproyecto{} (figura~\ref{fig:cadena-valor}) distingue las actividades primarias, que crean directamente el servicio, de las actividades de apoyo, que las sostienen.

Las \textbf{actividades primarias} comprenden: la recepción de contenido generado por creadores y de datos de uso (logística de entrada); la operación de la plataforma, incluida la disponibilidad de campañas, la sincronización de eventos en tiempo real y la administración del \textit{marketplace} (operaciones); la distribución continua de nuevas versiones y del contenido adquirido (logística de salida); la captación y retención de usuarios y creadores (marketing y ventas); y el soporte, la moderación y la incorporación de mejoras basadas en la retroalimentación de la comunidad (servicios).

Las \textbf{actividades de apoyo} son la infraestructura de la empresa (planificación, gestión financiera, cumplimiento normativo, monetización y alianzas), la administración de recursos humanos, el desarrollo de tecnologías (métricas de uso, rendimiento, seguridad e investigación de tecnologías emergentes) y las compras (equipamiento y alojamiento en infraestructura propia, con servicios en la nube como opción de escalado, herramientas de desarrollo, licencias, monitoreo y procesamiento de pagos).

Dado el carácter digital del servicio, el margen se origina sobre todo en la calidad de las operaciones y del desarrollo tecnológico, y en la capacidad de atraer contenido de la comunidad, más que en una cadena logística física.

\begin{landscape}
\begin{figure}[p]
  \centering
  \includegraphics[width=\linewidth,height=0.78\textheight,keepaspectratio]{2.2.2 Analisis de Cadena de Valor.png}
  \caption{Cadena de valor de \nombreproyecto{}. Elaboración propia.}
  \label{fig:cadena-valor}
\end{figure}
\end{landscape}

\subsubsection{Factores críticos de éxito}\label{sec:fce}
Los factores que determinarán el éxito del negocio son:
\begin{itemize}
  \item \textbf{Facilidad de uso:} condición para la adopción, en especial por parte de usuarios que se inician en el género.
  \item \textbf{Estabilidad y rendimiento:} la plataforma se utiliza durante la sesión de juego, por lo que no tolera interrupciones.
  \item \textbf{Crecimiento sostenido de la comunidad:} el valor de la plataforma aumenta con la cantidad de usuarios activos.
  \item \textbf{Generación constante de contenido:} un catálogo renovado sostiene el interés y alimenta el \textit{marketplace}.
  \item \textbf{Diferenciación frente a los VTT tradicionales:} el posicionamiento como complemento de la partida presencial es la base de la propuesta.
  \item \textbf{Escalabilidad tecnológica:} la infraestructura debe acompañar el crecimiento de usuarios sin deterioro del servicio.
  \item \textbf{Capacidad de retener usuarios activos:} el uso recurrente es la base de la conversión a planes premium.
  \item \textbf{Modelo económico sostenible:} los ingresos por suscripciones, comisiones y servicios adicionales deben cubrir los costos operativos.
\end{itemize}

\subsubsection{Fortalezas y debilidades del negocio}\label{sec:fortalezas-debilidades-sector}
Este apartado identifica las fortalezas y debilidades propias del sector de las herramientas digitales para juegos de rol de mesa, es decir, las del contexto en que el negocio competirá. Las fortalezas y debilidades específicas de la organización y del producto se analizan en el capítulo~3.

\paragraph{Fortalezas.}
\begin{itemize}
  \item \textbf{Expansión sostenida de la demanda:} la industria de los juegos de rol de mesa crece de forma acelerada en el mundo, impulsada por la exposición en medios masivos, por las plataformas de \textit{streaming} (como los canales de \textit{Actual Play} en Twitch y YouTube) y por videojuegos de rol de alcance mundial. Esto asegura un flujo constante de nuevos usuarios.
  \item \textbf{Comunidad altamente colaborativa:} el sector cuenta con una base de usuarios proactiva, muy inclinada a crear y distribuir contenido propio (campañas, mapas, reglas caseras o \textit{homebrew}). Esta dinámica facilita la viabilidad de los modelos de negocio basados en plataformas colaborativas y de intermediación, como los esquemas de usuario a usuario (C2C).
  \item \textbf{Madurez de la infraestructura tecnológica:} el avance de los servicios en la nube (\textit{cloud computing}) y la adopción masiva de dispositivos móviles permiten a las empresas del sector ofrecer soluciones escalables a nivel global con costos de infraestructura física prácticamente nulos.
  \item \textbf{Escalabilidad internacional sin fronteras físicas:} al tratarse de productos y servicios puramente digitales (SaaS), el sector puede cruzar fronteras de manera inmediata y alcanzar mercados internacionales sin las fricciones de la logística tradicional.
\end{itemize}

\paragraph{Debilidades.}
\begin{itemize}
  \item \textbf{Alta fragmentación de las herramientas:} el ecosistema de software para TTRPG carece de una integración estandarizada. Los usuarios deben emplear múltiples aplicaciones desconectadas entre sí para resolver distintas necesidades, como gestores de personajes, simuladores de dados, plataformas de mapas y aplicaciones de notas.
  \item \textbf{Elevada barrera técnica y curva de aprendizaje:} las soluciones dominantes, orientadas a la virtualización completa de la mesa (\textit{Virtual Tabletop} o VTT), tienen interfaces complejas que exigen a los directores de juego (DM) una inversión de tiempo considerable para dominarlas, lo que aleja a los usuarios principiantes.
  \item \textbf{Baja disposición al pago (alta sensibilidad al precio):} el nicho incluye un segmento importante de usuarios habituados a resolver sus necesidades con métodos tradicionales (papel y lápiz) o con herramientas digitales gratuitas y de uso general, como procesadores de texto en la nube o aplicaciones de mensajería. Esta resistencia cultural dificulta la retención en esquemas de suscripción de alto costo.
  \item \textbf{Dependencia crítica de propiedad intelectual de terceros:} la gran mayoría de las herramientas del sector estructura sus funcionalidades en torno a sistemas de reglas que pertenecen a grandes editoriales. Esto expone a la industria del software a un riesgo constante ante cambios en las políticas de licencias abiertas (como la \textit{Open Gaming License} de Wizards of the Coast) y limita la soberanía comercial de las empresas desarrolladoras.
\end{itemize}
